% ==============================================================================
% Scientific Slides — Academic Seminar Template (45–60 min)
% LaTeX Beamer, 16:9, Madrid theme with custom colors
% Compile: pdflatex seminar.tex
% ==============================================================================

\documentclass[aspectratio=169,11pt]{beamer}

\usepackage[utf8]{inputenc}
\usepackage[T1]{fontenc}
\usetheme{Madrid}
\usecolortheme{dolphin}
\setbeamertemplate{navigation symbols}{}

% Section divider pages
\AtBeginSection[]{
  \begin{frame}
    \vfill
    \centering
    \begin{beamercolorbox}[sep=8pt,center,shadow=true,rounded=true]{title}
      \usebeamerfont{title}\insertsectionhead\par%
    \end{beamercolorbox}
    \vfill
  \end{frame}
}

\usepackage{graphicx}
\graphicspath{{./figures/}}
\usepackage{amsmath, amssymb}
\usepackage{booktabs}
\usepackage{multirow}
\usepackage[style=authoryear,maxcitenames=2,backend=biber]{biblatex}
\addbibresource{references.bib}
\renewcommand*{\bibfont}{\tiny}

% Custom colors
\definecolor{darkblue}{RGB}{0,75,135}
\definecolor{lightblue}{RGB}{100,150,200}
\setbeamercolor{structure}{fg=darkblue}
\setbeamercolor{title}{fg=darkblue}
\setbeamercolor{frametitle}{fg=darkblue}

% Title
\title[Short Title]{Full Title of Your Research:\\Comprehensive and Descriptive}
\subtitle{Research Seminar Presentation}
\author[Your Name]{Your Name, PhD Candidate\\
  Advisor: Prof.\ Advisor Name}
\institute[University]{
  Department of Your Field\\
  University Name\\
  \texttt{yourname@university.edu}
}
\date{\today}

\begin{document}

% --- TITLE ---
\begin{frame}[plain]
  \titlepage
\end{frame}

\begin{frame}{Outline}
  \tableofcontents
\end{frame}

% ============================
% INTRODUCTION (8–10 min)
% ============================
\section{Introduction}

\begin{frame}{Motivation}
  \begin{columns}[T]
    \begin{column}{0.5\textwidth}
      \textbf{The Big Picture:}
      \begin{itemize}
        \item Why this research area matters
        \item Real-world impact and applications
        \item Current challenges in the field
      \end{itemize}
    \end{column}
    \begin{column}{0.5\textwidth}
      % \includegraphics[width=\textwidth]{motivation.pdf}
      \framebox[0.9\textwidth][c]{[Motivating Figure]}
    \end{column}
  \end{columns}
  \vspace{0.5cm}
  \begin{block}{Central Question}
    How can we address this challenge using novel approach X?
  \end{block}
\end{frame}

\begin{frame}{Prior Work}
  \textbf{Historical Development:}
  \begin{itemize}
    \item Early work established foundation \cite{foundational}
    \item Key advances in 2000s \cite{advance2005}
    \item Recent developments \cite{recent2022}
  \end{itemize}
  \vspace{0.5cm}
  \textbf{Current State:}
  \begin{enumerate}
    \item Known: X affects Y
    \item Evidence: mechanism involves Z
    \item Unknown: questions remain about W
  \end{enumerate}
\end{frame}

\begin{frame}{Knowledge Gap}
  \begin{columns}[c]
    \begin{column}{0.6\textwidth}
      \textbf{What We Know:}
      \begin{itemize}
        \item Established finding 1
        \item Replicated result 2
        \item General consensus 3
      \end{itemize}
      \vspace{0.3cm}
      \textbf{What Remains Unknown:}
      \begin{itemize}
        \item \alert{Gap 1:} Critical unknown
        \item \alert{Gap 2:} Methodological limitation
        \item \alert{Gap 3:} Unexplored context
      \end{itemize}
    \end{column}
    \begin{column}{0.4\textwidth}
      \begin{alertblock}{The Problem}
        Existing approaches fail to account for X, limiting understanding of Y.
      \end{alertblock}
    \end{column}
  \end{columns}
\end{frame}

\begin{frame}{Research Objectives}
  \begin{exampleblock}{Overall Goal}
    To investigate how X influences Y under conditions Z.
  \end{exampleblock}
  \vspace{0.5cm}
  \textbf{Specific Aims:}
  \begin{enumerate}
    \item \textbf{Aim 1:} Characterize X–Y relationship
    \item \textbf{Aim 2:} Identify mechanism W
    \item \textbf{Aim 3:} Test generalizability to context Z
  \end{enumerate}
\end{frame}

% ============================
% METHODS (8–10 min)
% ============================
\section{Methods}

\begin{frame}{Study Design}
  \begin{figure}
    \centering
    % \includegraphics[width=0.85\textwidth]{study_design.pdf}
    \framebox[0.8\textwidth][c]{[Study Design Schematic]}
    \caption{Multi-phase experimental design}
  \end{figure}
  \begin{itemize}
    \item \textbf{Phase 1:} Observational ($n = 150$)
    \item \textbf{Phase 2:} Controlled experiment ($n = 80$)
    \item \textbf{Phase 3:} Validation ($n = 120$)
  \end{itemize}
\end{frame}

\begin{frame}{Analysis Plan}
  \textbf{Primary Analyses:}
  \begin{itemize}
    \item \textbf{Aim 1:} Linear regression: $Y = \beta_0 + \beta_1 X + \epsilon$
    \item \textbf{Aim 2:} Mediation analysis (bootstrap, 5000 iterations)
    \item \textbf{Aim 3:} Mixed-effects model for context effects
  \end{itemize}
  \vspace{0.5cm}
  \begin{block}{Software}
    R 4.x (lme4, lavaan); Python 3.x (scikit-learn)
  \end{block}
\end{frame}

% ============================
% RESULTS (18–22 min)
% ============================
\section{Results}

\begin{frame}{Aim 1: X Predicts Y}
  \begin{columns}[c]
    \begin{column}{0.6\textwidth}
      \begin{figure}
        % \includegraphics[width=\textwidth]{aim1.pdf}
        \framebox[0.9\textwidth][c]{[Regression Plot]}
        \caption{$R^2 = 0.29$, $p < .001$}
      \end{figure}
    \end{column}
    \begin{column}{0.4\textwidth}
      \begin{table}
        \centering\tiny
        \begin{tabular}{lccc}
          \toprule
          \textbf{Predictor} & $\beta$ & SE & $p$ \\
          \midrule
          Intercept & 12.45 & 3.21 & <.001 \\
          X & 0.54 & 0.08 & <.001 \\
          \bottomrule
        \end{tabular}
      \end{table}
      \begin{block}{Finding}
        X significantly predicts Y, controlling for demographics.
      \end{block}
    \end{column}
  \end{columns}
\end{frame}

\begin{frame}{Aim 2: Mediation by W}
  \begin{figure}
    % \includegraphics[width=0.75\textwidth]{mediation.pdf}
    \framebox[0.7\textwidth][c]{[Mediation Diagram]}
    \caption{W mediates X $\to$ Y}
  \end{figure}
  \begin{itemize}
    \item \textbf{Direct effect:} $c' = 0.31$, $p = .021$
    \item \textbf{Indirect effect:} $ab = 0.23$, 95\% CI [0.14, 0.35]
    \item \textbf{Proportion mediated:} 43\%
  \end{itemize}
\end{frame}

\begin{frame}{Aim 3: Generalization}
  \begin{columns}[T]
    \begin{column}{0.5\textwidth}
      % \includegraphics[width=\textwidth]{context1.pdf}
      \framebox[0.9\textwidth][c]{[Context 1]}
    \end{column}
    \begin{column}{0.5\textwidth}
      % \includegraphics[width=\textwidth]{context2.pdf}
      \framebox[0.9\textwidth][c]{[Context 2]}
    \end{column}
  \end{columns}
  \vspace{0.5cm}
  \begin{itemize}
    \item Main effect of X: $\beta = 0.51$, $p < .001$
    \item Context $\times$ X interaction: ns ($p = .23$)
    \item \alert{Effect generalizes across contexts}
  \end{itemize}
\end{frame}

% ============================
% DISCUSSION (10–12 min)
% ============================
\section{Discussion}

\begin{frame}{Key Results and Interpretation}
  \begin{exampleblock}{Summary}
    \begin{enumerate}
      \item X predicts Y ($\beta = 0.54$, 29\% variance)
      \item W mediates 43\% of effect
      \item Generalizes across contexts
    \end{enumerate}
  \end{exampleblock}
  \vspace{0.5cm}
  \textbf{Relation to prior work:}
  \begin{itemize}
    \item Consistent with \cite{jones2020}
    \item Extends beyond prior scope
    \item Resolves contradictions in literature
  \end{itemize}
\end{frame}

\begin{frame}{Limitations and Future Directions}
  \textbf{Limitations:}
  \begin{enumerate}
    \item Cross-sectional design $\to$ future longitudinal
    \item University sample $\to$ community samples
    \item Self-report measures $\to$ objective measures
  \end{enumerate}
  \vspace{0.5cm}
  \begin{block}{Next Steps}
    \begin{itemize}
      \item Replicate in independent sample
      \item Develop intervention targeting W
      \item Investigate neural mechanisms
    \end{itemize}
  \end{block}
\end{frame}

% ============================
% CONCLUSION
% ============================
\section{Conclusion}

\begin{frame}[plain]
  \begin{center}
    {\LARGE \textbf{Thank You}}\\[1cm]
    {\Large Questions \& Discussion}\\[1.5cm]
    Your Name\\
    Department, University\\
    \texttt{yourname@university.edu}\\[1cm]
    {\footnotesize Funded by Grant Agency (\#12345)}
  \end{center}
\end{frame}

% --- BACKUP ---
\appendix

\begin{frame}{Backup: Full Regression Table}
  \framebox[\textwidth][c]{[Complete regression results with all covariates]}
\end{frame}

\begin{frame}[allowframebreaks]{References}
  \printbibliography
\end{frame}

\end{document}
